两方案均从1月1日6000\,kWh按各自正式策略连续运行；1月历史不足时采用周期预测，随后逐步积累同一预测流程的历史误差。2月1日期初储电量分别为6699.9484和1790.8119\,kWh，评价期均为2月1日至12月31日334日、48096槽，末日实际储电量分别为1200.0000和1200.0000\,kWh。方案4-2的混合整数规划单次限时5\,s、目标间隙0.2\%，固定模式购电修正最多120次迭代。其全年365次规划均取得可行解，混合整数规划与固定模式修正的累计耗时分别为1341.24、600.33\,s，实际最大相对间隙为2.7128\%。这些间隙只度量单次代理规划，不构成全年真实费用的最优性界。方案4-3采用HiGHS求解1460次线性规划，均获得可行解，全年运行与归档耗时127.19\,s。

图\ref{fig:q4-price}按同一48096个午夜预测槽比较价格模型，8维模型RMSE为0.06195元/kWh，10维模型为0.05614元/kWh。价格误差的下降反映预测精度改善，经济影响还需经过购电和储能决策。

评价期账单见表\ref{tab:q4-annual}。方案4-2与4-3总费用分别为14008609.54元、13539823.62元，差额为468785.91元（3.35\%）。方案4-3下调退款为83052.64元，紧急购电量为46082.4095\,kWh；与方案4-2相比紧急购电量减少58.87\%。两者的非空方向反转分别为1757与6679次，这些指标仅计评价期内动作，是运行强度的描述，不能直接换算电池寿命。
